% Options for packages loaded elsewhere
% Options for packages loaded elsewhere
\PassOptionsToPackage{unicode}{hyperref}
\PassOptionsToPackage{hyphens}{url}
\PassOptionsToPackage{dvipsnames,svgnames,x11names}{xcolor}
%
\documentclass[
  letterpaper,
  DIV=11,
  numbers=noendperiod]{scrartcl}
\usepackage{xcolor}
\usepackage{amsmath,amssymb}
\setcounter{secnumdepth}{-\maxdimen} % remove section numbering
\usepackage{iftex}
\ifPDFTeX
  \usepackage[T1]{fontenc}
  \usepackage[utf8]{inputenc}
  \usepackage{textcomp} % provide euro and other symbols
\else % if luatex or xetex
  \usepackage{unicode-math} % this also loads fontspec
  \defaultfontfeatures{Scale=MatchLowercase}
  \defaultfontfeatures[\rmfamily]{Ligatures=TeX,Scale=1}
\fi
\usepackage{lmodern}
\ifPDFTeX\else
  % xetex/luatex font selection
\fi
% Use upquote if available, for straight quotes in verbatim environments
\IfFileExists{upquote.sty}{\usepackage{upquote}}{}
\IfFileExists{microtype.sty}{% use microtype if available
  \usepackage[]{microtype}
  \UseMicrotypeSet[protrusion]{basicmath} % disable protrusion for tt fonts
}{}
\makeatletter
\@ifundefined{KOMAClassName}{% if non-KOMA class
  \IfFileExists{parskip.sty}{%
    \usepackage{parskip}
  }{% else
    \setlength{\parindent}{0pt}
    \setlength{\parskip}{6pt plus 2pt minus 1pt}}
}{% if KOMA class
  \KOMAoptions{parskip=half}}
\makeatother
% Make \paragraph and \subparagraph free-standing
\makeatletter
\ifx\paragraph\undefined\else
  \let\oldparagraph\paragraph
  \renewcommand{\paragraph}{
    \@ifstar
      \xxxParagraphStar
      \xxxParagraphNoStar
  }
  \newcommand{\xxxParagraphStar}[1]{\oldparagraph*{#1}\mbox{}}
  \newcommand{\xxxParagraphNoStar}[1]{\oldparagraph{#1}\mbox{}}
\fi
\ifx\subparagraph\undefined\else
  \let\oldsubparagraph\subparagraph
  \renewcommand{\subparagraph}{
    \@ifstar
      \xxxSubParagraphStar
      \xxxSubParagraphNoStar
  }
  \newcommand{\xxxSubParagraphStar}[1]{\oldsubparagraph*{#1}\mbox{}}
  \newcommand{\xxxSubParagraphNoStar}[1]{\oldsubparagraph{#1}\mbox{}}
\fi
\makeatother


\usepackage{longtable,booktabs,array}
\newcounter{none} % for unnumbered tables
\usepackage{calc} % for calculating minipage widths
% Correct order of tables after \paragraph or \subparagraph
\usepackage{etoolbox}
\makeatletter
\patchcmd\longtable{\par}{\if@noskipsec\mbox{}\fi\par}{}{}
\makeatother
% Allow footnotes in longtable head/foot
\IfFileExists{footnotehyper.sty}{\usepackage{footnotehyper}}{\usepackage{footnote}}
\makesavenoteenv{longtable}
\usepackage{graphicx}
\makeatletter
\newsavebox\pandoc@box
\newcommand*\pandocbounded[1]{% scales image to fit in text height/width
  \sbox\pandoc@box{#1}%
  \Gscale@div\@tempa{\textheight}{\dimexpr\ht\pandoc@box+\dp\pandoc@box\relax}%
  \Gscale@div\@tempb{\linewidth}{\wd\pandoc@box}%
  \ifdim\@tempb\p@<\@tempa\p@\let\@tempa\@tempb\fi% select the smaller of both
  \ifdim\@tempa\p@<\p@\scalebox{\@tempa}{\usebox\pandoc@box}%
  \else\usebox{\pandoc@box}%
  \fi%
}
% Set default figure placement to htbp
\def\fps@figure{htbp}
\makeatother


% definitions for citeproc citations
\NewDocumentCommand\citeproctext{}{}
\NewDocumentCommand\citeproc{mm}{%
  \begingroup\def\citeproctext{#2}\cite{#1}\endgroup}
\makeatletter
 % allow citations to break across lines
 \let\@cite@ofmt\@firstofone
 % avoid brackets around text for \cite:
 \def\@biblabel#1{}
 \def\@cite#1#2{{#1\if@tempswa , #2\fi}}
\makeatother
\newlength{\cslhangindent}
\setlength{\cslhangindent}{1.5em}
\newlength{\csllabelwidth}
\setlength{\csllabelwidth}{3em}
\newenvironment{CSLReferences}[2] % #1 hanging-indent, #2 entry-spacing
 {\begin{list}{}{%
  \setlength{\itemindent}{0pt}
  \setlength{\leftmargin}{0pt}
  \setlength{\parsep}{0pt}
  % turn on hanging indent if param 1 is 1
  \ifodd #1
   \setlength{\leftmargin}{\cslhangindent}
   \setlength{\itemindent}{-1\cslhangindent}
  \fi
  % set entry spacing
  \setlength{\itemsep}{#2\baselineskip}}}
 {\end{list}}
\usepackage{calc}
\newcommand{\CSLBlock}[1]{\hfill\break\parbox[t]{\linewidth}{\strut\ignorespaces#1\strut}}
\newcommand{\CSLLeftMargin}[1]{\parbox[t]{\csllabelwidth}{\strut#1\strut}}
\newcommand{\CSLRightInline}[1]{\parbox[t]{\linewidth - \csllabelwidth}{\strut#1\strut}}
\newcommand{\CSLIndent}[1]{\hspace{\cslhangindent}#1}



\setlength{\emergencystretch}{3em} % prevent overfull lines

\providecommand{\tightlist}{%
  \setlength{\itemsep}{0pt}\setlength{\parskip}{0pt}}



 


% Define chapter counter if document class doesn't have one (e.g. article)
\ifdefined\thechapter\else
  \newcounter{chapter}
  \renewcommand*{\thechapter}{\arabic{chapter}}
\fi

\usepackage{setspace}
\onehalfspacing

\usepackage[pagewise]{lineno}
\linenumbers

\usepackage{xcolor}
\usepackage{tikz}
\usepackage{afterpage}
\usepackage{eso-pic}
\definecolor{thesiscream}{RGB}{255, 252, 240}
\definecolor{thesissand}{RGB}{232, 212, 162}
\pagecolor{thesiscream}
\usepackage{booktabs}
\usepackage{longtable}
\usepackage{array}
\usepackage{multirow}
\usepackage{wrapfig}
\usepackage{float}
\usepackage{colortbl}
\usepackage{pdflscape}
\usepackage{tabularx}
\usepackage{xltabular}
\usepackage{threeparttable}
\usepackage{threeparttablex}
\usepackage[normalem]{ulem}
\usepackage{makecell}
\usepackage{xcolor}
\KOMAoption{captions}{tableheading}
\makeatletter
\@ifpackageloaded{caption}{}{\usepackage{caption}}
\AtBeginDocument{%
\ifdefined\contentsname
  \renewcommand*\contentsname{Table of contents}
\else
  \newcommand\contentsname{Table of contents}
\fi
\ifdefined\listfigurename
  \renewcommand*\listfigurename{List of Figures}
\else
  \newcommand\listfigurename{List of Figures}
\fi
\ifdefined\listtablename
  \renewcommand*\listtablename{List of Tables}
\else
  \newcommand\listtablename{List of Tables}
\fi
\ifdefined\figurename
  \renewcommand*\figurename{Figure}
\else
  \newcommand\figurename{Figure}
\fi
\ifdefined\tablename
  \renewcommand*\tablename{Table}
\else
  \newcommand\tablename{Table}
\fi
}
\@ifpackageloaded{float}{}{\usepackage{float}}
\floatstyle{ruled}
\@ifundefined{c@chapter}{\newfloat{codelisting}{h}{lop}}{\newfloat{codelisting}{h}{lop}[chapter]}
\floatname{codelisting}{Listing}
\newcommand*\listoflistings{\listof{codelisting}{List of Listings}}
\makeatother
\makeatletter
\makeatother
\makeatletter
\@ifpackageloaded{caption}{}{\usepackage{caption}}
\@ifpackageloaded{subcaption}{}{\usepackage{subcaption}}
\makeatother
\usepackage{bookmark}
\IfFileExists{xurl.sty}{\usepackage{xurl}}{} % add URL line breaks if available
\urlstyle{same}
% fallback for those not using the hyperref driver hyperxmp:
\makeatletter
\@ifundefined{xmpquote}{\newcommand{\xmpquote}[1]{#1}}{}
\makeatother
\hypersetup{
  pdftitle={Stone Piles as Kōura Habitat},
  pdfauthor={Olivier V. Raven; Ian A. K. Kusabs; Robin Holmes; Francis J. Burdon; Deniz Özkundakci},
  pdfkeywords={\xmpquote{Ameiurus nebulosus}, \xmpquote{Artificial
reef}, \xmpquote{Freshwater conservation}, \xmpquote{Habitat
enhancement}, \xmpquote{Kōura}, \xmpquote{Occupancy
modelling}, \xmpquote{Paranephrops planifrons}, \xmpquote{Restoration
ecology}},
  colorlinks=true,
  linkcolor={blue},
  filecolor={Maroon},
  citecolor={Blue},
  urlcolor={Blue},
  pdfcreator={LaTeX via pandoc}}


\title{Stone Piles as Kōura Habitat}
\usepackage{etoolbox}
\makeatletter
\providecommand{\subtitle}[1]{% add subtitle to \maketitle
  \apptocmd{\@title}{\par {\large #1 \par}}{}{}
}
\makeatother
\subtitle{Stone pile refuges in an Aotearoa New Zealand lake: early
benefits for threatened crayfish but not its invasive predator}
\author{Olivier V. Raven \and Ian A. K. Kusabs \and Robin
Holmes \and Francis J. Burdon \and Deniz Özkundakci}
\date{2026-09-25}
\begin{document}
\maketitle


\clearpage
\providecommand{\chaptermark}[1]{}
\ifcsname c@chapter\endcsname
\setcounter{chapter}{5}
\setcounter{figure}{0}
\setcounter{table}{0}
\setcounter{section}{0}
\addcontentsline{toc}{chapter}{\protect\numberline{5}Stone Piles as K\={o}ura Habitat}
\chaptermark{Stone Piles as K\={o}ura Habitat}
\fi
\providecommand{\thechapter}{5}
\thispagestyle{empty}
\begin{tikzpicture}[remember picture, overlay]
  \begin{scope}
    \clip ([yshift=-7cm]current page.north west)
          rectangle ([yshift=4.5cm]current page.south east);
    \node[inner sep=0pt, anchor=north]
      at ([yshift=-7cm]current page.north)
      {\includegraphics[height=18.2cm]{images/ch5_cover.jpg}};
  \end{scope}
  \fill[black]
    (current page.north west) rectangle ([yshift=-7cm]current page.north east);
  \node[color=thesissand, align=left,
        font=\fontsize{13}{16}\selectfont\sffamily, anchor=west]
    at ([xshift=1.5cm, yshift=-1.5cm]current page.north west)
    {Chapter \thechapter};
  \node[color=white, text width=0.85\paperwidth, align=left,
        font=\fontsize{22}{28}\selectfont\bfseries\sffamily, anchor=west]
    at ([xshift=1.5cm, yshift=-4cm]current page.north west)
    {Stone Piles as K\={o}ura Habitat};
  \node[color=white, text width=0.85\paperwidth, align=left,
        font=\fontsize{10}{13}\selectfont\itshape\sffamily, anchor=north west]
    at ([xshift=1.5cm, yshift=-5.0cm]current page.north west)
    {Stone pile refuges in an Aotearoa New Zealand lake: early benefits for threatened crayfish but not its invasive predator};
  \fill[black]
    (current page.south west) rectangle ([yshift=4.5cm]current page.south east);
  \node[color=white, text width=0.85\paperwidth, align=left,
        font=\fontsize{9}{13}\selectfont\sffamily, anchor=west]
    at ([xshift=1.5cm, yshift=3.2cm]current page.south west)
    {Olivier V.\ Raven, Ian A.\ K.\ Kusabs, Robin Holmes, Francis J.\ Burdon \& Deniz Özkundakci};
  \node[color=thesissand, text width=0.85\paperwidth, align=left,
        font=\fontsize{9}{13}\selectfont\itshape\sffamily, anchor=west]
    at ([xshift=1.5cm, yshift=1.5cm]current page.south west)
    {Restoration Ecology--- Under review};
  \node[color=white!80, anchor=south west,
        font=\fontsize{7}{9}\selectfont\itshape\sffamily]
    at ([xshift=1.5cm, yshift=5cm]current page.south west)
    {Stone pile construction at Lake Rotoiti.};
\end{tikzpicture}
\null
\clearpage

\subsection*{Abstract}\label{abstract}
\addcontentsline{toc}{subsection}{Abstract}

\textbf{Introduction}: Kōura (\emph{Paranephrops planifrons}), Aotearoa
New Zealand's endemic freshwater crayfish, have declined by up to 91\%
in Lake Rotoiti (Te Arawa Lakes, Rotorua), driven by predation from the
invasive brown bullhead catfish (\emph{Ameiurus nebulosus}) and thermal
stratification that compresses suitable habitat into the shallow
littoral zone, stressors that cannot be removed at the lake scale.

\textbf{Objectives}: We tested whether stone pile structures could
provide kōura with structural refuge from catfish predation and
established a transferable construction methodology for lake managers.

\textbf{Methods}: Kōura and catfish were monitored at 20 locations (five
stone piles, five control, and ten reference sites) over 18 months using
a quasi before-after control-impact (quasi-BACI) design. Kōura occupancy
was modelled within a Bayesian framework incorporating detection
probabilities from baited traps and fyke nets, alongside a
simulation-based power analysis.

\textbf{Results}: Stone pile construction produced no statistically
significant change in kōura abundance, biomass, or occupancy relative to
controls, but effects consistently trended positive across all models.
Kōura use was strongly seasonal, falling in winter as kōura withdrew to
deeper water at stone pile and control sites, while reference sites
remained stable; catfish showed no comparable attraction. The study was
underpowered, given high pre-construction kōura occupancy (81--89\%) and
only five site pairs.

\textbf{Conclusions}: Stone piles show early, consistent evidence of
selectively benefiting kōura without benefiting its invasive predator,
though the study lacked power to statistically confirm this trend.

\textbf{Implications for Practice}: Stone piles are most valuable when
habitat compression and predation pressure peak together in summer and
autumn, conditions expected to intensify under climate warming. This
study establishes a tested, scalable construction and consenting
methodology that lake managers can deploy as an interim refuge-based
tool while catfish eradication remains unachievable at the lake scale.
We recommend continued monitoring beyond three years, given the slow
trajectories typical of crayfish colonisation, and trialling
corridor-style structures linking refuge to deeper water rather than
isolated piles, to improve winter resilience and support scaling toward
population-level recovery. As catfish showed no comparable benefit from
the structures in this trail, managers can pursue stone pile refuges for
kōura without the risk of also benefiting this invasive predator.

\subsection{Introduction}\label{introduction}

Freshwater biodiversity is declining faster than any other biome, with
one quarter of freshwater fauna now threatened with extinction
(\citeproc{ref-Sayer2025}{Sayer et al. 2025}). Among the most endangered
are freshwater crayfish, with 32\% of species assessed as threatened
globally due to habitat loss, invasive species, and declining water
quality (\citeproc{ref-Richman2015}{Richman et al. 2015}). As keystone
species and ecosystem engineers, crayfish play crucial roles in nutrient
cycling, bioturbation, and food web dynamics, and their loss can have
cascading consequences for freshwater ecosystems
(\citeproc{ref-Collier1997}{Collier et al. 1997};
\citeproc{ref-Parkyn2001}{Parkyn et al. 2001}). In Aotearoa New Zealand,
these pressures converge on kōura (\emph{Paranephrops planifrons} White,
1842), an endemic freshwater crayfish now facing rapid population
decline in lakes that have long supported some of the largest customary
kōura fisheries in the country (\citeproc{ref-Hiroa1921}{Hiroa 1921};
\citeproc{ref-Kusabs2009}{Kusabs \& Quinn 2009}).

Kōura populations in Lake Rotoiti, located in the Rotorua Te Arawa Lakes
region of the North Island of Aotearoa New Zealand (hereafter Aotearoa),
have declined by up to 91\% over the past two decades
(\citeproc{ref-Kusabs2026}{Kusabs et al. 2026}). This collapse has
virtually ended customary harvest practices for Te Arawa iwi (Māori
tribal group) who are hunga tiaki (traditional custodians) of the lake
(\citeproc{ref-Kusabs2015a}{Kusabs, Hicks, et al. 2015};
\citeproc{ref-Kusabs2026}{Kusabs et al. 2026}). The decline in kōura
numbers is driven by multiple interacting stressors that compound over
time in different parts of the lake (\citeproc{ref-Lee2025}{Lee et al.
2025}; \citeproc{ref-Kusabs2026}{Kusabs et al. 2026}). The invasive
brown bullhead catfish (\emph{Ameiurus nebulosus} (Lesueur, 1819);
hereafter catfish), confirmed in Lake Rotoiti in 2016
(\citeproc{ref-Francis2019}{Francis 2019};
\citeproc{ref-Grayling2020}{Grayling 2020}), preys directly on kōura and
suppresses their foraging through non-consumptive `fear' effects
(\citeproc{ref-Raven2026a}{Raven, MacNeil, et al. 2026}). Predation
pressure from catfish is compounded by additional stressors that
constrain where in the lake kōura can seek refuge from this invasive
predator (\citeproc{ref-Kusabs2026}{Kusabs et al. 2026}).

During the day, kōura shelter in complex rocky habitat and in deeper
water to avoid predation (\citeproc{ref-Jowett2008}{Jowett et al. 2008};
\citeproc{ref-Kusabs2015b}{Kusabs, Quinn, et al. 2015}), moving out of
both at night to forage in the productive littoral zone
(\citeproc{ref-Devcich1979}{Devcich 1979}). The importance of deep-water
refuge is shown in Lake Taupō, where kōura persist at greater depth even
in the presence of catfish because the lake remains oxygenated below the
thermocline (\citeproc{ref-Kusabs2015c}{Kusabs \& Taiaroa 2015};
\citeproc{ref-Verburg2019}{Verburg \& Albert 2019}). In Lake Rotoiti,
this deep refuge becomes anoxic during summer stratification,
compressing the zone of liveable habitat for kōura into the shallow
littoral margins (\citeproc{ref-Kusabs2026}{Kusabs et al. 2026}),
precisely where catfish are most active during summer
(\citeproc{ref-Dedual2002}{Dedual 2002};
\citeproc{ref-Blanchfield2026}{Blanchfield et al. 2026}). Dense beds of
invasive macrophytes further restrict kōura movement through this
compressed zone and intensify oxygen depletion during die-off events
(\citeproc{ref-Vincent1984}{Vincent et al. 1984};
\citeproc{ref-Hamilton2005}{Hamilton et al. 2005};
\citeproc{ref-Kusabs2009}{Kusabs \& Quinn 2009}). The dual pressures of
catfish predation and thermal stratification are expected to only
strengthen with predicted climate scenarios
(\citeproc{ref-Woolway2021}{Woolway et al. 2021};
\citeproc{ref-Lee2025}{Lee et al. 2025}), making kōura habitat
compression and predation pressure progressively more severe.

The compression of suitable habitat and subsequent concentration of
kōura into the littoral zone makes the quality of nearshore lacustrine
habitats increasingly important for remaining populations. In stream
environments, kōura are positively associated with structurally complex
habitats including bank undercuts, woody debris, leaf litter, and coarse
substrates (\citeproc{ref-Parkyn2004}{Parkyn \& Collier 2004};
\citeproc{ref-Jowett2008}{Jowett et al. 2008};
\citeproc{ref-Parkyn2009}{Parkyn et al. 2009}). In lake environments,
benthic substrate characteristics are a stronger predictor of kōura
distribution than nutrient enrichment or predator abundance
(\citeproc{ref-Kusabs2015b}{Kusabs, Quinn, et al. 2015}), and kōura are
positively associated with coarse substrates and riparian vegetation
(\citeproc{ref-Raven2026b}{Raven, Burdon, et al. 2026}). Coarse rocky
substrate provides refugia from predators and shelter during moulting
(\citeproc{ref-Streissl2002}{Streissl \& Hödl 2002}). Experimental
studies confirm that shelter availability directly reduces catfish
predation on crayfish (\citeproc{ref-Adams2007}{Adams 2007};
\citeproc{ref-Zhao2024}{Zhao et al. 2024}), and kōura increase refuge
use in the presence of catfish compared to native predators
(\citeproc{ref-Raven2026a}{Raven, MacNeil, et al. 2026}), emphasising
the importance of shelter availability for kōura populations. Shelter
availability is likely a limiting factor for kōura in Lake Rotoiti,
where the majority of the shoreline is sandy with dense beds of
introduced macrophytes and little rocky structure. Adding rocky
substrate to soft-bottomed lake areas has been shown to significantly
increase crayfish abundance (\citeproc{ref-Johnsen2008}{Johnsen \&
Taugbøl 2008}), and increasing structural complexity in streams improves
crayfish restocking success (\citeproc{ref-Huolila1997}{Huolila et al.
1997}). Deploying habitat enhancement structures that add structural
complexity in the littoral zone could potentially be a practical way to
restore kōura populations in Lake Rotoiti.

In this study habitat enhancement structures in the shape of stone piles
were deployed along the shoreline of Lake Rotoiti to investigate the
ability of the structures in providing additional habitat for kōura to
live and hide from predators. Kōura presence, abundance, and biomass
were monitored at stone pile, control, and reference sites over 18
months. We predicted that kōura occupancy, abundance and biomass would
increase more rapidly at the stone pile sites than at control sites, as
intense catfish predation pressure in the littoral zone was expected to
compel kōura actively toward available refuge structures. We further
predicted that kōura abundance would continue to increase progressively
after construction, at least within the first 18 months, and that
catfish would instead show seasonal patterns of shoreline use rather
than a construction-driven trend (\citeproc{ref-Barnes1996}{Barnes
1996}; \citeproc{ref-Dedual2002}{Dedual 2002};
\citeproc{ref-Blanchfield2026}{Blanchfield et al. 2026}). Finally, we
predicted that stone pile sites would progressively converge toward
kōura densities observed at natural rocky reference sites, and that
catfish would not exhibit a similar pattern.

\subsection{Methods}\label{methods}

\subsubsection{Study location and
construction}\label{study-location-and-construction}

Lake Rotoiti (-38.03, 176.42) is a warm monomictic mesotrophic lake
located in the Rotorua Te Arawa Lakes region of the North Island of
Aotearoa (34 km², maximum depth 128 m). The lake stratifies for
approximately nine months a year ({\textbf{?@fig-lake-conditions}}),
with pronounced thermal stratification compressing oxygenated habitat
into the shallow littoral zone during summer and autumn. Since 2008,
nutrient-rich inflows from Lake Rotorua have been diverted toward the
Kaituna River via the Ōhau Channel diversion wall to improve water
quality in Lake Rotoiti (\citeproc{ref-Hamilton2005}{Hamilton et al.
2005}; \citeproc{ref-Prentice2026}{Prentice et al. 2026}). Lake level is
regulated by a weir at the outlet of the Kaituna River, providing a
stable and controlled water level throughout the year. Catfish are
present throughout the lake but occur in particularly high
concentrations in Te Weta Bay on the Northwest side of the lake
(\citeproc{ref-Francis2019}{Francis 2019}). Five stone pile sites were
selected alongside the lake shoreline based on accessibility and
landowner approval (Figure~\ref{fig-sites}). At each stone pile site, a
paired control site was established 50 m along the shoreline, matched
for habitat characteristics e.g., water depth, substrate type and
vegetation cover. Five naturally occurring rocky shoreline locations,
with each two monitoring sites 50 m apart were selected as reference. In
total, 20 sites of each 150 m² (15 × 10 m) were monitored seasonally
over 18 months. Physical characteristics of all sites are summarised in
Table~\ref{tbl-site-overview}. Baseline monitoring was conducted in
February 2025 prior to stone pile deployment, which occurred the
following week. Stone piles were constructed from locally sourced quarry
rock (Manawahe Stone and Quarry, \textasciitilde25 km from the lake)
with a diameter ranging from 85 - 315 mm (mean = 173 mm). At each of the
five stone pile sites, three stone piles were arranged horizontally
along the shoreline by hand from the shore. Stone pile design was
informed by mesocosm experiments demonstrating kōura preference for
enclosed stone structures (unpublished data). Stone piles were
positioned approximately 4.3 m from shore at a mean depth of 0.49 m and
with a mean area of 4.1 m².

\begin{figure}

\centering{

\includegraphics[width=1\linewidth,height=\textheight,keepaspectratio]{images/Rotoiti_monitoring_sites.png}

}

\caption{\label{fig-sites}Map of Lake Rotoiti in Bay of Plenty region of
New Zealand. The monitoring sites are coloured by site type: stone pile
sites red, control site yellow, and reference sites blue. Ten close-up
maps show the placement of each of the sites.}

\end{figure}%

\begingroup\fontsize{8}{10}\selectfont

\begin{longtable}[t]{rllrrrrrr}

\caption{\label{tbl-site-overview}Physical characteristics of stone
pile, control, and reference monitoring sites at baseline. Sites 1--10
are stone pile (odd) and paired control (even) sites; sites 11--20 are
reference sites. Substrate index is a weighted composite of substrate
particle sizes (higher values indicate coarser substrate more suitable
for kōura). Slope is the gradient to the 20 m depth contour. Fetch is
maximum effective fetch. Vegetation values are visually estimated
percentage cover.}

\tabularnewline

\toprule
Site & Type & Habitat & Substrate index & Slope (m/m) & Fetch (km) & Riparian veg. (\%) & Overhang (\%) & Wood cover (\%)\\
\midrule
\cellcolor{gray!8}{1} & \cellcolor{gray!8}{Stone pile} & \cellcolor{gray!8}{Rocky} & \cellcolor{gray!8}{5.20} & \cellcolor{gray!8}{0.12} & \cellcolor{gray!8}{4.76} & \cellcolor{gray!8}{0} & \cellcolor{gray!8}{10} & \cellcolor{gray!8}{0}\\
2 & Control & Rocky & 4.40 & 0.11 & 4.79 & 70 & 0 & 0\\
\cellcolor{gray!8}{3} & \cellcolor{gray!8}{Stone pile} & \cellcolor{gray!8}{Sandy} & \cellcolor{gray!8}{3.00} & \cellcolor{gray!8}{0.01} & \cellcolor{gray!8}{1.71} & \cellcolor{gray!8}{100} & \cellcolor{gray!8}{0} & \cellcolor{gray!8}{0}\\
4 & Control & Sandy & 3.00 & 0.01 & 1.74 & 80 & 5 & 0\\
\cellcolor{gray!8}{5} & \cellcolor{gray!8}{Stone pile} & \cellcolor{gray!8}{Sandy} & \cellcolor{gray!8}{2.40} & \cellcolor{gray!8}{0.01} & \cellcolor{gray!8}{0.61} & \cellcolor{gray!8}{10} & \cellcolor{gray!8}{0} & \cellcolor{gray!8}{0}\\
\addlinespace
6 & Control & Sandy & 3.00 & 0.01 & 0.55 & 0 & 0 & 0\\
\cellcolor{gray!8}{7} & \cellcolor{gray!8}{Stone pile} & \cellcolor{gray!8}{Sandy} & \cellcolor{gray!8}{3.00} & \cellcolor{gray!8}{0.06} & \cellcolor{gray!8}{2.50} & \cellcolor{gray!8}{100} & \cellcolor{gray!8}{0} & \cellcolor{gray!8}{0}\\
8 & Control & Sandy & 3.00 & 0.05 & 2.22 & 30 & 0 & 0\\
\cellcolor{gray!8}{9} & \cellcolor{gray!8}{Stone pile} & \cellcolor{gray!8}{Sandy} & \cellcolor{gray!8}{3.35} & \cellcolor{gray!8}{0.08} & \cellcolor{gray!8}{3.59} & \cellcolor{gray!8}{0} & \cellcolor{gray!8}{0} & \cellcolor{gray!8}{3}\\
10 & Control & Rocky & 5.00 & 0.09 & 3.74 & 20 & 0 & 0\\
\addlinespace
\cellcolor{gray!8}{11} & \cellcolor{gray!8}{Reference} & \cellcolor{gray!8}{Rocky} & \cellcolor{gray!8}{6.60} & \cellcolor{gray!8}{0.16} & \cellcolor{gray!8}{4.47} & \cellcolor{gray!8}{0} & \cellcolor{gray!8}{0} & \cellcolor{gray!8}{0}\\
12 & Reference & Rocky & 6.20 & 0.12 & 4.17 & 100 & 100 & 0\\
\cellcolor{gray!8}{13} & \cellcolor{gray!8}{Reference} & \cellcolor{gray!8}{Sandy} & \cellcolor{gray!8}{3.60} & \cellcolor{gray!8}{0.09} & \cellcolor{gray!8}{1.71} & \cellcolor{gray!8}{100} & \cellcolor{gray!8}{100} & \cellcolor{gray!8}{0}\\
14 & Reference & Sandy & 3.50 & 0.11 & 1.65 & 100 & 100 & 0\\
\cellcolor{gray!8}{15} & \cellcolor{gray!8}{Reference} & \cellcolor{gray!8}{Rocky} & \cellcolor{gray!8}{5.20} & \cellcolor{gray!8}{0.03} & \cellcolor{gray!8}{1.44} & \cellcolor{gray!8}{100} & \cellcolor{gray!8}{100} & \cellcolor{gray!8}{6}\\
\addlinespace
16 & Reference & Rocky & 5.20 & 0.04 & 1.40 & 100 & 100 & 2\\
\cellcolor{gray!8}{17} & \cellcolor{gray!8}{Reference} & \cellcolor{gray!8}{Rocky} & \cellcolor{gray!8}{5.00} & \cellcolor{gray!8}{0.05} & \cellcolor{gray!8}{3.47} & \cellcolor{gray!8}{100} & \cellcolor{gray!8}{100} & \cellcolor{gray!8}{30}\\
18 & Reference & Rocky & 6.20 & 0.05 & 3.46 & 100 & 100 & 7\\
\cellcolor{gray!8}{19} & \cellcolor{gray!8}{Reference} & \cellcolor{gray!8}{Rocky} & \cellcolor{gray!8}{5.70} & \cellcolor{gray!8}{0.11} & \cellcolor{gray!8}{2.63} & \cellcolor{gray!8}{100} & \cellcolor{gray!8}{100} & \cellcolor{gray!8}{4}\\
20 & Reference & Rocky & 5.70 & 0.13 & 2.87 & 100 & 100 & 0\\
\bottomrule

\end{longtable}

\endgroup{}

\textsubscript{Source:
\href{https://OlivierRaven.github.io/Reefs_Rotoiti/index.qmd.html}{Article
Notebook}}

\subsubsection{Monitoring}\label{monitoring}

Physical, chemical, and biological parameters were recorded at each site
during each monitoring round. Static site characteristics including
substrate composition, slope to depth contour, riparian vegetation,
overhanging tree cover, and wood cover were visually estimated during
baseline monitoring only, as these were not expected to change
substantially over the monitoring period. Substrate composition was
visually estimated by a single observer walking around each site and
using an underwater viewing chamber, categorising bedrock, boulders
(\textgreater256 mm), cobble (64--256 mm), gravel (2--64 mm), sand
(0.063--2 mm), mud (\textless0.063 mm), and organic matter
(\citeproc{ref-Wentworth1922}{Wentworth 1922}). A substrate index, that
indicates suitability for kōura, was calculated from these estimates
(\citeproc{ref-Jowett1990}{Jowett \& Richardson 1990};
\citeproc{ref-Jowett2008}{Jowett et al. 2008}):

\[
\begin{aligned}
\text{Substrate index} &= 0.08 \times \text{bedrock} + 0.07 \times \text{boulders} \\
&\quad + 0.06 \times \text{cobble} + 0.04 \times \text{gravel} + 0.03 \times \text{sand} \\
&\quad + 0.02 \times \text{organic matter} + 0.01 \times \text{mud}
\end{aligned}
\] Coarser substrates were assigned higher weights as they provide
greater interstitial pore space and stable structures used by kōura
(\citeproc{ref-Usio2000}{Usio \& Townsend 2000};
\citeproc{ref-Nystrom2006}{Nyström et al. 2006}).

After stone pile construction, measurements of the size and dimensions
of each structure were taken. The post-construction substrate index for
stone pile sites was recalculated using each site's combined stone-pile
footprint area, expressed as a percentage of the 150 m² site area and
treated as 80\% cobble and 20\% boulders displacing an equivalent area
of sand, to reflect their contribution to substrate complexity. Slope
was assessed as the horizontal distance from the shoreline to the 20 m
depth contour, derived from bathymetric maps and measured in a projected
coordinate system (NZTM2000; EPSG:2193).

Seasonal variable parameters were recorded at every monitoring round and
included aquatic vegetation cover (emergent and submerged macrophytes).
Chemical parameters included water temperature (°C), dissolved oxygen
(mg/L), and conductivity (μS/cm) were measured using a YSI ProSolo meter
(Pro2030, YSI Inc., Yellow Springs, OH, USA). pH was measured using a
handheld pH meter (EC-PCTestr35, Eutech Instruments).

Kōura and fish were sampled every monitoring round using two
single-winged, double-throated fyke nets (600 mm high D-shaped entrance
hoop) and two collapsible bait traps (480 × 250 mm; two 50 mm entrance
openings) baited with catfish flesh (\citeproc{ref-Joy2013}{Joy et al.
2013}; \citeproc{ref-MacNeil2025}{MacNeil et al. 2025}). Fyke nets were
placed perpendicular to the shoreline outside of the stone pile
structures or at 10 m intervals, with bait traps positioned between
them. All equipment was deployed simultaneously at each site for a 24-h
period. Captured kōura were measured for orbital carapace length (OCL)
using a digital calliper (0.01 mm precision; 150 mm) and for wet mass
using a digital scale (± 1 g; 5 kg capacity, Anko Brand, Kmart Group,
Mulgrave, VIC, Australia). Catfish and other fish species, were recorded
for abundance, measured for total length, and weighed for wet mass.
Catch per unit effort (CPUE) and biomass per unit effort (BPUE) were
calculated for all kōura and fish based on the number of nets and traps
deployed at each site per monitoring round, standardised by soak time
(\citeproc{ref-Kusabs2015b}{Kusabs, Quinn, et al. 2015}). Kōura weights
were unavailable for 36 individuals because of scale failure and were
therefore estimated from a log₁₀--log₁₀ length--weight regression fitted
to individuals with both measurements (\emph{n} = 840), following
standard approaches for weight length relationships in aquatic organisms
{\textbf{?@fig-length-weight-ch5}} (\citeproc{ref-Froese2006}{Froese
2006}). Predicted values were back-transformed with a lognormal bias
correction factor \(10^{0.5 \sigma^2}\) (\citeproc{ref-Quinn2002}{Quinn
\& Keough 2002}).

\[
\log_{10}(\text{Weight [g]}) = -2.67 + 2.684 \times \log_{10}(\text{OCL length [mm]})
\]

\subsubsection{Data analysis}\label{data-analysis}

Data analysis was a three-part process, the first assessed baseline
equivalence between stone pile, control, and reference sites before
construction, the second assessed the impact of stone pile construction
and after-construction trajectory on kōura and catfish occupancy,
abundance, and biomass, and the third assessed whether stone pile sites
converged toward reference sites and if reference sites remained stable
over time. All analyses were conducted in R (version 4.6.1; R Core Team
(\citeproc{ref-RCoreTeam2026}{2026})) using the glmmTMB and emmeans
packages (\citeproc{ref-Brooks2017}{Brooks et al. 2017};
\citeproc{ref-Lenth2017}{Lenth \& Piaskowski 2017}).

Covariate selection for all models followed a consistent procedure.
Corrected Akaike Information Criterion (AICc) values for each base model
were compared against models with additional covariates (temperature,
season, substrate index, maximum effective fetch, or wave energy proxy)
added at a time, and a covariate was retained only when it reduced AICc
by more than two units relative to base model
(\citeproc{ref-Burnham2002}{Burnham \& Anderson 2002}). Robustness and
sensitivity for all models were assessed by evaluating uniform
distribution, dispersion and outlier rate using the DHARMa package
(\citeproc{ref-Hartig2016}{Hartig 2016}). Across the full set of models,
the Benjamini--Hochberg correction method was applied to control the
false discovery rate within each of the three analysis components.

Site variation before stone pile construction was assessed using paired
Wilcoxon signed-rank tests to compare stone pile sites with their paired
control sites, and unpaired Wilcoxon rank-sum tests to compare stone
pile sites with reference sites. This was done for response variables
(water quality, substrate, and kōura and catfish CPUE and BPUE) to
confirm that the stone pile and control sites were equivalent before
construction, and to characterise the difference between stone pile and
reference sites. To further assess baseline equivalence, generalised
linear models (GLMs) were fitted to test whether kōura and catfish CPUE
and BPUE differed between site type (stone pile / control) and between
the two reference sites (reference 1 / reference 2) using a Tweedie
family with a log link. Following the covariate selection procedure,
temperature was included in the catfish models comparing reference
sites, as the base model failed to converge and temperature was required
for numerical stability.

The second component of the analysis tested construction impact relative
to control sites using a quasi-before-after, control-impact (quasi-BACI)
design (\citeproc{ref-Underwood1994}{Underwood 1994}), and separately
tested after-construction trajectory of kōura and catfish populations.
The design is quasi-BACI because only a single pre-construction
monitoring round was possible. The quasi-BACI design compared one
pre-construction monitoring round with six post-construction monitoring
rounds using generalised linear mixed effects models (GLMMs) to test the
effect of stone pile construction on kōura and catfish occupancy,
abundance, and biomass. Time period (before / after) and site type
(stone pile / control) were included as an interaction term, monitoring
round as a covariate and site nested within control-impact pair as a
random effect. Occupancy was modelled using a binomial GLMM, and
abundance and biomass using a Tweedie GLMM with a log link; fetch was
retained as a covariate in the catfish model. The post-construction
trajectory models used the same response variables, families, covariate,
and random effect structure, but with site type and monitoring round
included as an interaction term in place of the before / after
comparison.

The third component of the analysis assessed whether stone pile sites
converged toward reference sites over time, and whether reference sites
themselves remained stable. Kōura and catfish abundance and biomass were
tested using GLMMs with site type (stone pile / reference) and
monitoring round included as an interaction term, and site number as a
random effect, using a Tweedie family with log link; temperature was
retained as a covariate in the catfish model. Reference site stability
was tested in GLMMs with site type (reference 1 / reference 2) and
monitoring round as an interaction term, and site nested within pair as
a random effect, again using Tweedie family with log link; season was
retained as a covariate in the kōura model and temperature was retained
as a covariate in the catfish model.

To complement the GLMM abundance analysis, kōura site-level occupancy
was estimated using a two-method occupancy model fitted within a
Bayesian framework. Because bait traps and fyke nets provide independent
and simultaneous detection opportunities at each site, true site
occupancy could be separated from imperfect and method-specific
detection probabilities. A custom Stan observation model was implemented
using the brms package (\citeproc{ref-Burkner2017}{Bürkner 2017},
\citeproc{ref-Burkner2018}{2018}) using cmdstanr
(\citeproc{ref-Gabry2025}{Gabry et al. 2025}) to estimate true occupancy
and detection probabilities for baited traps and fyke nets via logit
links. Standardised water temperature was included as the covariate for
detection probability. Three occupancy sub-models were fitted. A full
model across all monitoring rounds served as baseline comparison. A
summer-only model and a warm-season model (summer and autumn combined)
were fitted to focus on the period of highest thermal stratification,
when kōura are known to disproportionately occupy nearshore littoral
habitats (\citeproc{ref-Kusabs2026}{Kusabs et al. 2026}) and when
catfish predation pressure is highest in the areas where the stone piles
are located (\citeproc{ref-Blanchfield2026}{Blanchfield et al. 2026}).
The warm-season model maximises ecological relevance, while retaining a
tractable formula with season as an additive covariate. All models were
run with four chains of 4000 iterations (1500 warmup) and convergence
was assessed using R̂ values (threshold \textless{} 1.01) and bulk
effective sample sizes (threshold \textgreater{} 400)
(\citeproc{ref-Vehtari2021}{Vehtari et al. 2021}). Divergent transitions
were rare in all models (≤ 9 of 10,000 post-warmup draws) and considered
negligible given the otherwise strong convergence diagnostics.

A simulation-based power analysis was conducted to quantify the sample
size required for statistical detection of the BACI effect under the
observed effect size and variance structure, and to put the
non-significant results in perspective relative to study design rather
than true absence of effect. For each analysis the occupancy and
abundance, datasets were simulated across a grid of scenarios varying in
number of site pairs (3-30) and after-period monitoring rounds (2-12).
For the occupancy BACI, simulated binary detection data were generated
using values from the warm-season occupancy model posterior means, and
each simulated dataset was analysed using a logistic GLMM (lme4; Bates
et al. (\citeproc{ref-Bates2015}{2015})). For the CPUE BACI, simulated
abundance data were generated using fixed effects and variance
components from the kōura abundance BACI model, using a log-normal
approximation to the Tweedie distribution, and each dataset was analysed
using a linear mixed model fitted to a log-transformed CPUE. Statistical
power was estimated as the proportion of simulations (300 per scenario)
in which the model's 95\% confidence interval for the BACI coefficient
excluded zero in the positive direction.

\subsection{Results}\label{results}

Kōura presence was detected at stone pile sites within three months of
construction, with occupancy and abundance trending consistently higher
at stone pile sites than paired control sites throughout the 18-month
post-construction period, though this difference did not reach
statistical significance. Results across all models and response
variables are summarised in Figure~\ref{fig-forest-plot} and
{\textbf{?@tbl-models-summary-full-supp}}.

\begin{figure}

\centering{

\includegraphics[width=1\linewidth,height=\textheight,keepaspectratio]{outputs/fig-forest-plot.png}

}

\caption{\label{fig-forest-plot}Effect sizes (\(\beta\) ± 95\% CI) for
the key term of interest across all models, separated by species
(columns) and comparison type (rows). Filled points indicate significant
effects (\emph{p} \textless{} 0.05); open points are non-significant.
Each row has its own x-axis scale, effect magnitudes are not directly
comparable across rows because each comparison tests a different
contrast (step-change, per-round slope, or simple group difference), and
occupancy is on the log-odds scale while abundance and biomass models
use the Tweedie log link. Models with degenerate confidence intervals
are excluded.}

\end{figure}%

\subsubsection{Baseline equivalence}\label{baseline-equivalence}

Pre-construction, kōura abundance and biomass did not differ
significantly between stone pile and paired control sites
({\textbf{?@tbl-baseline-combined-supp}}), confirming baseline
equivalences (\emph{p} = 0.783 and \emph{p} = 0.571). Kōura abundance
(\emph{p} = 0.968) and biomass (\emph{p} = 0.968) did not differ between
reference sites. Catfish abundance (\emph{p} = 0.004) and biomass
(\emph{p} = 0.002) likewise did not differ between stone pile and
control sites pre-construction. However, catfish abundance and biomass
were significantly higher at Reference 1 than Reference 2 (\emph{p} =
0.004 and \emph{p} = 0.002, respectively), despite both sites having
comparable rocky substrate. This suggests catfish distribution varies
considerably between nearby sites with similar habitat characteristics.

\subsubsection{Construction impact (BACI) and post-construction
trajectory}\label{construction-impact-baci-and-post-construction-trajectory}

Stone pile construction produced no detectable step-change in kōura
abundance, biomass, or occupancy relative to paired control sites over
the 18-month monitoring period following deployment (all
\emph{p}\textsubscript{adj} = 0.858) (Figure~\ref{fig-koura-effect}). A
matched-season comparison, comparing Summer 2025 (before construction)
with Summer 2026 (one year after construction), produced the same
conclusion (kōura abundance: \emph{p} = 0.582). Post-construction
trajectories did not diverge significantly between stone pile and
control sites (all \emph{p}\textsubscript{adj} = 0.215). Kōura abundance
declined significantly across post-construction rounds at control and
stone pile sites (\(\beta\) = -0.41, \emph{p} = 0.004), with the
sharpest decline observed in Winter 2026 (Figure~\ref{fig-cpue-trends}).
Stone pile construction produced no detectable step-change or trajectory
divergence in catfish abundance, biomass, or occupancy relative to
control sites across any model or time period
(Figure~\ref{fig-forest-plot}).

\begin{figure}

\centering{

\includegraphics[width=1\linewidth,height=\textheight,keepaspectratio]{outputs/fig-koura-effect.png}

}

\caption{\label{fig-koura-effect}Kōura CPUE at individual site pairs
across all monitoring rounds. Each panel shows one stone pile / control
pair (sites 1--10) or reference pair (sites 11--20). Colour and shape
indicate site type. Vertical dotted line marks stone pile construction.
y-axis on log(1+x) scale.}

\end{figure}%

\begin{figure}

\centering{

\includegraphics[width=1\linewidth,height=\textheight,keepaspectratio]{outputs/fig-cpue-trends.png}

}

\caption{\label{fig-cpue-trends}Mean CPUE per monitoring round (± 95\%
CI, log scale) by site type for kōura (top) and catfish (bottom).
Vertical dashed line distinguishes the pre-construction period (Before)
from post-construction monitoring (After).}

\end{figure}%

Kōura occupancy was high across all site types before construction,
81--89\% in summer ({\textbf{?@tbl-psi-warm-supp}}). After construction,
occupancy increased marginally at all sites, 89--94\% in summer, thereby
limiting our ability to detect a treatment effect. The BACI contrast to
compare how much stone pile sites changed relative to control sites was
slightly positive, but the credible intervals spanned zero with wide
uncertainty (Figure~\ref{fig-occ-warm-contrasts}), indicating no
detectable additional increase at stone pile sites. However, this result
likely reflects a measured ceiling rather than biological response.
After-treatment contrast showed no significant difference between site
types, but the direction of the effect was consistent with early habitat
use: stone pile sites showed a positive contrast relative to control
sites, while control sites showed a negative contrast relative to
reference sites, and stone pile sites were near zero relative to
reference sites (Figure~\ref{fig-occ-warm-contrasts}).

\begin{figure}

\centering{

\includegraphics[width=0.8\linewidth,height=\textheight,keepaspectratio]{outputs/fig-occ-warm-contrasts.png}

}

\caption{\label{fig-occ-warm-contrasts}Posterior contrasts in kōura
occupancy from the warm-season model. BACI contrast (stone pile minus
control, before vs after construction). After-treatment contrasts
between all site-type pairs. Points show posterior means with 95\%
credible intervals; contrasts overlapping zero indicate no detectable
difference.}

\end{figure}%

The non-significant BACI results should be interpreted in the context of
statistical power. For occupancy, power was near zero across all
simulated scenarios (maximum 8\%), reflecting the ceiling effect leaving
no statistical room for a BACI signal
({\textbf{?@fig-power-curves-combined}}). Post-construction power
analysis showed that achieving 80\% power with the current five site
pairs would require an effect size of approximately 1.8 log-units,
almost 6 times the current observed effect
(Figure~\ref{fig-trajectory-projection}). If the estimated positive
trajectory of +0.24 log-units per seasonal round continues, the effect
size may reach this detection threshold within approximately 6
additional monitoring rounds (\textasciitilde1.5 years).

\begin{figure}

\centering{

\includegraphics[width=0.8\linewidth,height=\textheight,keepaspectratio]{outputs/fig-trajectory-projection.png}

}

\caption{\label{fig-trajectory-projection}Projected kōura CPUE BACI
effect size over additional monitoring rounds. The solid line shows the
linear projection based on the estimated post-construction trajectory;
shaded ribbon shows 95\% uncertainty interval. Dashed horizontal line
marks the detection threshold for 80\% power with five site pairs.}

\end{figure}%

\subsubsection{Convergence toward reference sites and reference
stability}\label{convergence-toward-reference-sites-and-reference-stability}

Stone pile sites maintained a consistently lower kōura abundance than
reference sites throughout the post-construction period, with no
significant change in this gap over time (\(\beta\) = -0.09, \emph{p} =
0.429; {\textbf{?@fig-convergence-plot}}). Kōura abundance declined
across monitoring rounds at both stone pile and reference sites
(\(\beta\) = -0.09, \emph{p} = 0.107). The Winter 2026 monitoring round
showed a pronounced decline at stone pile sites (CPUE ≈ 0.7) while
reference sites maintained substantially higher abundance (CPUE ≈ 1.7;
Figure~\ref{fig-cpue-trends}), a within-round contrast not captured by
the linear trajectory model. Water temperature at Winter 2026 was
approximately 1.5°C lower than Winter 2025 but consistent across site
types, ruling out temperature as a site-specific driver. Visual
inspection of stone pile surfaces during Winter 2026 monitoring revealed
substantially greater filamentous algal accumulation compared to earlier
rounds ({\textbf{?@fig-stone-piles}}). Catfish abundance and biomass
showed no significant convergence or divergence between stone pile and
reference sites over time (Figure~\ref{fig-forest-plot}).

\subsection{Discussion}\label{discussion}

Stone pile structures were deployed in the littoral zone of Lake Rotoiti
to test whether physical refuge could reduce predation pressure on kōura
by invasive catfish. Stone piles produced no statistically significant
increase in kōura abundance, biomass, or occupancy relative to paired
control sites, but the direction of effect was consistently positive
across every response variable and model, and occupancy at stone pile
sites converged with occupancy at natural rocky reference sites by the
first post-construction summer. Catfish showed no comparable response to
stone pile presence in any model, at any site type, or in any season,
indicating that stone piles selectively benefit kōura without providing
equivalent refuge or foraging opportunity for this invasive predator. A
pronounced decline in kōura catch rates at stone pile and control sites
in Winter 2026, evidently driven by seasonal redistribution to deeper
water rather than habitat abandonment, showed that structural refuge
value is itself seasonal and concentrated in the warmer months when
catfish predation pressure and habitat compression peak. Together, these
results and a formal power analysis indicate that the absence of
statistical significance most plausibly reflects limited study duration
and sample size rather than genuine ineffectiveness of the intervention.

\subsubsection{Construction effect and statistical
power}\label{construction-effect-and-statistical-power}

Kōura catch rates were highly variable across all site types throughout
the monitoring period, with wide confidence intervals, likely reflecting
genuine biological variability in kōura distribution. This variability
may reflect seasonal activity patterns, detection probability changes
across monitoring rounds, and natural heterogeneity in littoral zone
use, rather than stable site associations. Despite this variability, the
Bayesian occupancy model revealed a consistent positive directional
trend, with stone pile sites showing higher occupancy than control
sites, and occupancy comparable to natural rocky reference sites across
both summer and autumn monitoring rounds. Whether this reflects active
refuge-seeking driven by catfish predation pressure or gradual
colonisation and population expansion into newly available habitats
cannot be determined from the current data, but the direction of the
response was consistent with the purpose of the structures, yet no
statistically significant effect was detected within the 18-month
monitoring period.

The absence of a significant effect is most likely explained by a
combination of monitoring duration, statistical power, and system
dynamics rather than ineffectiveness of the structures. Habitat
establishment for freshwater crayfish can be inherently slow. Johnsen \&
Taugbøl (\citeproc{ref-Johnsen2008}{2008}) reported that crayfish
abundance at a stone addition site took 13 years to reach reference
levels, and monitoring periods shorter than two years are often
insufficient to detect significant responses to habitat modification
(\citeproc{ref-Testa2011}{Testa et al. 2011};
\citeproc{ref-Ruh2012}{Ruhí et al. 2012};
\citeproc{ref-Pilotto2018}{Pilotto et al. 2018}). Following
reconstruction of the Lake Rotorua lakefront, which is connected to Lake
Rotoiti, kōura showed no meaningful recolonisation after 28 months
(\citeproc{ref-Kusabs2023}{Kusabs 2023}). They attributed slow
recolonisation to persistent catfish predation pressure and distance
from source population, consistent with time scales observed here. The
power analysis indicates that five site pairs are insufficient to
reliably detect effects of the magnitude observed here within 18 months,
establishing a minimum design threshold: future stone pile monitoring
programmes should plan for at least three years of monitoring and expect
high pre-construction occupancy, leaving little statistical room for a
measurable treatment effect. The simultaneous elevation of kōura CPUE
across all site types in Autumn 2025 further complicated detection of an
early treatment effect in the BACI contrast.

Features of the system itself may also contribute to the weak response.
Kōura abundance is highest in lakes without catfish
(\citeproc{ref-Raven2026b}{Raven, Burdon, et al. 2026}), suggesting
predation pressure may suppress kōura populations well below the
carrying capacity of existing habitat. If so, structural habitat
availability may not currently be the primary limiting factor, and a
weak or delayed response to habitat addition would be expected. Equally,
if kōura are already at carrying capacity of existing rocky habitat,
stone piles add new capacity that will only be filled gradually as the
population expands, which is unlikely to be detectable within 18 months.
In either case, a delayed rather than absent population response would
be expected. Isolated stone pile structures may also be difficult for
kōura to locate, since distance between isolated habitat patches is a
key determinant of colonisation success
(\citeproc{ref-Hanski1998}{Hanski 1998}). Each of these explanations
point toward a delayed, rather than absent, population response, and
continued monitoring would be the appropriate next step rather than a
reassessment of the intervention itself.

\subsubsection{Seasonal dynamics: the Winter 2026
decline}\label{seasonal-dynamics-the-winter-2026-decline}

Kōura CPUE at stone pile and control sites fell sharply in Winter 2026,
returning to levels close to the pre-deployment baseline from Summer
2025, while reference sites maintained substantially higher abundance.
This pattern is consistent with known seasonal migration behaviour
across lake depth profiles and kōura activity rates during cooler water
temperatures (\citeproc{ref-Devcich1979}{Devcich 1979}). When lake
thermal stratification breaks down in winter, the hypolimnion
re-oxygenates and deep-water habitat becomes accessible, allowing kōura
to redistribute along the depth gradient
(\citeproc{ref-Devcich1979}{Devcich 1979}). Reduced activity at low
water temperatures independently lowers the probability of trap capture
(\citeproc{ref-Devcich1979}{Devcich 1979}), and although Winter 2026 was
1.5 °C lower than Winter 2025, this difference was consistent across all
site types and therefore cannot explain why reference sites were
substantially more resilient to the winter decline than the stone pile
sites.

Natural rocky reference sites in this study represent the ecological
target for stone pile deployment. Reference sites had established kōura
populations, connectivity to deeper waters, and extensive refuge
habitat, allowing kōura to maintain littoral presence through winter or
return rapidly from deeper refuges as conditions fluctuated. Stone pile
sites are isolated structures with less connectivity to deep and
extensive refuge habitat, meaning kōura that withdrew in winter had no
quick pathway back to the structures. This difference likely explains
why stone pile sites maintained consistently lower kōura abundance than
reference sites throughout the monitoring period, with no significant
convergence toward reference levels. However, kōura occupancy at stone
pile sites was comparable to that at reference sites after construction,
suggesting that kōura found and used the structures at similar frequency
to natural rocky habitat but at lower densities. This distinction
between occupancy and abundance likely reflects the early stage of
colonisation and the small structural footprint of stone piles relative
to the extensive natural rocky habitat at reference sites. These
differences are expected to narrow as biofilm succession matures and
kōura accumulate at the structures over time.

Extensive filamentous algae growth was observed on the stone pile
structures in Winter 2026 but not to the same degree in Winter 2025,
consistent with the seasonal pattern typical for thermally stratified
lakes, where winter mixing raises nutrient availability while cooler
water suppresses invertebrate grazing, allowing algal mats to accumulate
on hard surfaces (\citeproc{ref-Wetzel2001}{Wetzel 2001};
\citeproc{ref-Bonacina2023}{Bonacina et al. 2023}). The freshly quarried
piles were free from biofilm at deployment and by Winter 2026 supported
a diverse meiofaunal assemblage, indicating active succession rather
than habitat degradation; algal cover would be expected to decline as
grazer activity resumes in spring and summer. Despite the near-zero trap
catches in Winter 2026, night spotlighting at two stone pile sites in
August 2026 directly observed kōura sheltering within the structures,
confirming that low winter catch rates reflect reduced detectability
rather than abandonment. Together, these results indicate that stone
pile refuge value is inherently seasonal --- greatest in summer and
autumn, when thermal stratification compresses habitat into the littoral
zone and catfish predation peaks, and lowest in winter, when kōura
redistribute to deeper water regardless of structure type.

\subsubsection{Catfish: no evidence of attraction to stone
piles}\label{catfish-no-evidence-of-attraction-to-stone-piles}

Catfish showed their seasonal patterns in littoral zone use, but these
were unrelated to stone pile presence. Stone pile construction produced
no evidence of catfish attraction or accumulation. Even if catfish were
to use stone piles to some degree, the ecological significance of this
would be limited given the extensive availability of existing suitable
habitat for catfish in Lake Rotoiti, particularly invasive macrophyte
beds (\citeproc{ref-Barnes2003}{Barnes \& Hicks 2003}). The single-site
increase in catfish CPUE during Summer 2026 was mirrored by an increase
at reference sites and was consistent with lake-wide seasonal
redistribution into warm littoral waters rather than any
stone-pile-specific event (\citeproc{ref-Blanchfield2026}{Blanchfield et
al. 2026}). Overall catfish catches remained low throughout the study
relative to targeted removal netting in Te Weta Bay, where highest
catfish densities are concentrated (\citeproc{ref-Francis2019}{Francis
2019}). In autumn, catfish shift to mid-water movement across the lake
basin remaining above 15 m depth
(\citeproc{ref-Blanchfield2026}{Blanchfield et al. 2026}), reducing
predation pressure on kōura in the littoral zone. This absence of
attraction is an important management result in its own right: habitat
enhancement can be targeted at a threatened native species without
inadvertently improving conditions for its invasive predator.

\subsubsection{Restoration design and future
monitoring}\label{restoration-design-and-future-monitoring}

Predator-prey interactions are shaped by habitat structure, which offers
refuge to prey, influences how easily predators locate and capture prey,
and alters the costs and benefits of hunting for a predator
(\citeproc{ref-Lennox2025}{Lennox et al. 2025}). This framework applies
directly to catfish-invaded lakes, where habitat addition offers a
practical lever on that interaction, and rocky environments have been
identified as natural multi-threat refugia, capable of simultaneously
buffering predation, thermal extremes, and invasive species pressure
across a wide range of taxa and ecosystems
(\citeproc{ref-Selwood2020}{Selwood \& Zimmer 2020}). In Lake Rotoiti,
catfish and thermal stratification are the two primary drivers of kōura
decline, and neither can be addressed at the lake scale. Catfish
eradication is not currently achievable, and targeted catfish removal in
high density areas reduces local pressure without eliminating the
lake-wide population (\citeproc{ref-Dedual2019}{Dedual 2019}). Thermal
stratification, meanwhile is expected to intensify under climate warming
(\citeproc{ref-Woolway2021}{Woolway et al. 2021};
\citeproc{ref-Lee2025}{Lee et al. 2025}). Where these stressors cannot
be removed, habitat addition becomes one of the few available management
tools, and refugia become correspondingly more valuable
(\citeproc{ref-Selwood2020}{Selwood \& Zimmer 2020}).

The stone piles deployed here were isolated structures, separated from
each other and from deeper water by soft sandy substrate. Isolated
habitat patches support fewer species and are more vulnerable to local
extinction than connected habitats (\citeproc{ref-Hanski1998}{Hanski
1998}), and the Winter 2026 decline is consistent with this limitation:
stone piles lacked the connectivity to deep refuge that allowed
reference sites to remain stable. Orienting future structures
perpendicular to the shoreline, extending through macrophyte beds and
into deeper water, could create corridors linking the littoral zone with
deeper refugia, allowing kōura to move more freely between depth ranges
as conditions change, as they do at reference sites. This design has not
been tested but is a promising direction for future deployments, with
the addition of potential benefit of suppressing macrophyte growth along
the corridor.

The dense algal fouling observed in Winter 2026 also raises a placement
question. In sheltered bays, stone piles are protected from wave damage,
but algal growth and sediment accumulation could progressively reduce
crevice accessibility and refugia value; in more exposed locations,
natural scouring could limit fouling, but conditions may be too
disturbed to provide effective shelter (\citeproc{ref-Parkyn2004}{Parkyn
\& Collier 2004}). This placement trade-off warrants direct
investigation in future deployments, and standardised photography of
consistent stone faces each monitoring visit would allow fouling
development to be tracked alongside changes in kōura use.

Finally, the five stone pile structures deployed here cannot be expected
to drive population-level recovery across a lake of this size on their
own. The spatial and temporal scale of the ecosystem responses is
typically greater than the scale of the interventions used to produce
them (\citeproc{ref-Minns1996}{Minns et al. 1996}), and scaling from
individual stone pile deployment to habitat addition at
population-relevant extents is a central challenge in restoration
ecology (\citeproc{ref-Ngwenya2025}{Ngwenya et al. 2025}).

Beyond the ecological results, this study establishes a practical
foundation for scaling habitat enhancement in these lakes: the
approvals, consents, and construction methodology needed to deploy stone
piles have now been tested, with evidence that doing so is unlikely to
have a counterproductive result by benefiting catfish. These outcomes
give lake managers a practical, tested option for future large-scale
deployment; continued monitoring for at least three years is recommended
to determine whether the Winter 2026 decline was a one-off event and
whether kōura abundance at stone pile sites continues to build in
subsequent spring and summer seasons.

Stone pile deployment alone cannot reverse the trajectory of kōura
decline in Lake Rotoiti, but it directly addresses the current
bottleneck: the shortage of structural refuge in a lake where soft
sediment dominates the shorelines and where the two primary stressors,
catfish predation and thermal habitat compression, cannot be controlled
at the lake scale. Both stressors are expected to intensify under
continued climate warming (\citeproc{ref-Woolway2021}{Woolway et al.
2021}; \citeproc{ref-Lee2025}{Lee et al. 2025}), narrowing the band of
habitat available to kōura and increasing, rather than diminishing, the
value of structural refugia over time. Establishing refugia before
populations decline further may be more effective than waiting until
source populations fall below critical densities for recolonisation
(\citeproc{ref-Morelli2016}{Morelli et al. 2016};
\citeproc{ref-Kusabs2026}{Kusabs et al. 2026}). The consistent positive
directional trend in occupancy, the rapid and clearly seasonal use of
the structures, and the developing biological community on the stone
pile surfaces together suggest the foundations for a population-level
response by kōura are being laid, even though catfish showed no
comparable benefit from the structures. Whether this translates into a
detectable population-level effect will depend on continued long-term
monitoring --- the logical next phase of this restoration programme.

\subsection{Acknowledgements}\label{acknowledgements}

We thank Te Arawa Lakes Trust and Te Komiti Whakahaere for the
opportunity to study kōura - a taonga (treasured) species. We are
grateful to Soweeta Fort-D'ath and William Anaru (Te Arawa Lakes Trust),
Tihini Grant (Ngāti Pikiao), and Andy Bruere (Bay of Plenty Regional
Council). We thank Tāheke Opatia Marae, Te Takinga Marae, and
landowners, Piripi Curtis, David McIlraith, Piki Thomas, Haddon Lock,
Barnett Vercoe, Richard Vercoe, and Caroline Newton, for granting access
to monitoring sites and facilitating fieldwork in Lake Rotoiti. We
further thank, Alex Simes, Cory O'Neill, Joe Butterworth, Martin
Sarkezi, Mael Marguet, Savio Dsouza, Ethan Morgan, Logan Kallam, Calum
MacNeil, and Matt Prentice for assisting with fieldwork and data
collection. Further thanks to Mael Marguet for suggesting the power
analysis.

\subsubsection{Funding statement}\label{funding-statement}

This research was supported by the Fish Futures programme funded through
a Ministry of Business, Innovation and Employment grant (CAWX2101) with
additional funding provided by the Bay of Plenty Regional Council under
the Toihuarewa Waimāori - Bay of Plenty Regional Council Chair in Lake
and Freshwater Science programme.

\subsubsection{Authors' contributions}\label{authors-contributions}

OVR, IAKK, RH, FJB, DÖ conceived and designed the research; OVR
performed the fieldwork; OVR analysed the data; IAKK obtained council
and iwi consent and provided local ecological knowledge; OVR drafted the
manuscript; OVR, IAKK, RH, FJB, DÖ revised the manuscript.

\subsubsection{Use of Artificial
Intelligence}\label{use-of-artificial-intelligence}

Use of Artificial Intelligence: Claude (Anthropic Sonnet 5.0) was used
to assist with manuscript editing and language refinement. All content
was critically reviewed, and the authors take full responsibility for
the final manuscript.

\subsubsection{Declarations}\label{declarations}

The corresponding author has declared that none of the authors has any
competing interests.

\subsubsection{Data availability
statement}\label{data-availability-statement}

The code and derived data supporting the findings of this study are
openly available on GitHub
(https://github.com/OlivierRaven/Stone\_Piles\_Rotoiti) and a rendered
version of the analysis is hosted at
https://olivierraven.github.io/Stone\_Piles\_Rotoiti/. The primary raw
dataset is archived on Zenodo (https://doi.org/\ldots). Peer review
documents, including the cover letter and reviewer responses, are
available in the manuscript repository in the interest of open and
transparent science.

\subsection{Supplementary information}\label{supplementary-information}

\emph{Supplementary figures are provided in the appendices.}

\protect\phantomsection\label{refs}
\begin{CSLReferences}{1}{1}
\bibitem[\citeproctext]{ref-Adams2007}
Adams SB (2007)
\href{https://doi.org/10.1674/0003-0031(2007)158\%5B85:daieoc\%5D2.0.co;2}{Direct
and indirect effects of channel catfish ({\emph{Ictalurus punctatus}})
on native crayfishes ({Cambaridae}) in experimental tanks}. The American
Midland Naturalist 158:85--96.

\bibitem[\citeproctext]{ref-Barnes1996}
Barnes GE (1996) The biology and general ecology of the brown bullhead
catfish (\emph{ameiurus nebulosus}) in {Lake Taupo}. PhD thesis,
University of Waikato, Hamilton, New Zealand

\bibitem[\citeproctext]{ref-Barnes2003}
Barnes GE, Hicks BJ (2003) Brown bullhead catfish ({\emph{Ameiurus
nebulosus}}) in {Lake Taupo}. Managing invasive freshwater fish in {New
Zealand}: Proceedings of a workshop hosted by {Department of
Conservation}, 10-12 may 2001, {Hamilton}. Department of Conservation

\bibitem[\citeproctext]{ref-Bates2015}
Bates D, Mächler M, Bolker B, Walker S (2015)
\href{https://doi.org/10.18637/jss.v067.i01}{Fitting linear
mixed-effects models using {lme4}}. Journal of Statistical Software
67:1--48.

\bibitem[\citeproctext]{ref-Blanchfield2026}
Blanchfield PJ, Rodrigues TH, MacNeil C, Stewart SD, Lee F, Hicks BJ,
Grayling S, Fort-D'Ath S, Prentice MJ, Baker C, et al. (2026)
\href{https://doi.org/10.1007/s10530-026-03877-5}{Invading brown
bullhead catfish (\emph{ameiurus nebulosus}) show thermal niche rigidity
and benthic niche plasticity: Implications for management}. Biological
Invasions 28.

\bibitem[\citeproctext]{ref-Bonacina2023}
Bonacina L, Fasano F, Mezzanotte V, Fornaroli R (2023)
\href{https://doi.org/10.1111/brv.12903}{Effects of water temperature on
freshwater macroinvertebrates: A systematic review}. Biological Reviews
98:191--221.

\bibitem[\citeproctext]{ref-Brooks2017}
Brooks M, Bolker B, Kristensen K, Maechler M, Magnusson A, Skaug H,
Nielsen A, Berg C, Bentham K van (2017)
\href{https://doi.org/10.32614/cran.package.glmmtmb}{glmmTMB:
Generalized linear mixed models using template model builder}. CRAN:
Contributed Packages.

\bibitem[\citeproctext]{ref-Burkner2018}
Bürkner P-C (2018) \href{https://doi.org/10.32614/rj-2018-017}{Advanced
{B}ayesian multilevel modeling with the {R} package brms}. The R Journal
10:395.

\bibitem[\citeproctext]{ref-Burkner2017}
Bürkner P-C (2017) \href{https://doi.org/10.18637/jss.v080.i01}{{brms}:
An {R} package for {B}ayesian multilevel models using {Stan}}. Journal
of Statistical Software 80.

\bibitem[\citeproctext]{ref-Burnham2002}
Burnham KP, Anderson DR (2002)
\href{https://doi.org/10.1007/978-0-387-22456-5_6}{Advanced issues and
deeper insights}. Pages 267--351. Model selection and multimodel
inference: A practical information-theoretic approach. Springer, New
York, NY

\bibitem[\citeproctext]{ref-Collier1997}
Collier KJ, Parkyn SM, Rabeni CF (1997) Koura: A keystone species. Water
and Atmosphere 5:18--20.

\bibitem[\citeproctext]{ref-Dedual2019}
Dedual M (2019)
\href{https://atlas.boprc.govt.nz/api/v1/edms/document/A3591759/content}{Summary
of the impacts and control methods of brown bullhead catfish
({\emph{Ameiurus nebulosus}} {Lesueur} 1815) in {New Zealand} and
overseas}. Bay of Plenty Regional Council

\bibitem[\citeproctext]{ref-Dedual2002}
Dedual M (2002)
\href{https://doi.org/10.1007/978-94-017-0771-8_15}{Vertical
distribution and movements of brown bullhead ({\emph{Ameiurus
nebulosus}} {Lesueur} 1819) in {Motuoapa Bay}, southern {Lake Taupo},
{New Zealand}}. Thorstad, EB, Fleming, IA, \& Næsje, TF (eds)
Hydrobiologia 483:129--135.

\bibitem[\citeproctext]{ref-Devcich1979}
Devcich AA (1979) An ecological study of {\emph{Paranephrops
planifrons}} {White} ({Decapoda} : {Parastacidae}) in {Lake Rotoiti},
{North Island}. PhD thesis, University of Waikato, Hamilton, New Zealand

\bibitem[\citeproctext]{ref-Francis2019}
Francis LB (2019)
\href{https://researchcommons.waikato.ac.nz/handle/10289/12746}{Evaluating
the effects of invasive brown bullhead catfish ({\emph{Ameiurus
nebulosus}}) on k{ō}ura (freshwater crayfish, {\emph{Paranephrops
planifrons}}) in {Lake Rotoiti}}. Master's thesis, University of Waikato

\bibitem[\citeproctext]{ref-Froese2006}
Froese R (2006)
\href{https://doi.org/10.1111/j.1439-0426.2006.00805.x}{Cube law,
condition factor and weight-length relationships: History, meta-analysis
and recommendations}. Journal of Applied Ichthyology 22:241--253.

\bibitem[\citeproctext]{ref-Gabry2025}
Gabry J, Češnovar R, Johnson A, Bronder S (2025)
\href{https://mc-stan.org/cmdstanr/}{Cmdstanr: R interface to
'CmdStan'}.

\bibitem[\citeproctext]{ref-Grayling2020}
Grayling S (2020)
\href{https://atlas.boprc.govt.nz/api/v1/edms/document/A3651696/content}{{Bay
of Plenty} regional pest management plan 2011-2016 annual report for
2019/2020}. Bay of Plenty Regional Council

\bibitem[\citeproctext]{ref-Hamilton2005}
Hamilton D, McBride C, Uraoka T (2005) {Lake Rotoiti} fieldwork and
modelling to support considerations of {Ohau} channel diversion from
{Lake Rotoiti}. NIWA \& University of Waikato

\bibitem[\citeproctext]{ref-Hanski1998}
Hanski I (1998) \href{https://doi.org/10.1038/23876}{Metapopulation
dynamics}. Nature 396:41--49.

\bibitem[\citeproctext]{ref-Hartig2016}
Hartig F (2016)
\href{https://doi.org/10.32614/cran.package.dharma}{{DHARMa}: Residual
diagnostics for hierarchical (multi-level / mixed) regression models}.
CRAN: Contributed Packages.

\bibitem[\citeproctext]{ref-Hiroa1921}
Hiroa TR (1921) M{ā}ori food-supplies of {Lake Rotorua}, with methods of
obtaining them, and usages and customs appertaining thereto.
Transactions and Proceedings of the Royal Society of New Zealand
1868-1961 53:433--451.

\bibitem[\citeproctext]{ref-Huolila1997}
Huolila M, Marjomäki TJ, Laukkanen E (1997)
\href{https://doi.org/10.1016/s0165-7836(97)00046-5}{The success of
crayfish stocking in a dredged river with and without artificial shelter
increase}. Fisheries Research 32:185--189.

\bibitem[\citeproctext]{ref-Johnsen2008}
Johnsen SI, Taugbøl T (2008) Add stones, get crayfish -- is it that
simple? Freshwater Crayfish 16:47--50.

\bibitem[\citeproctext]{ref-Jowett2008}
Jowett IG, Parkyn SM, Richardson J (2008)
\href{https://doi.org/10.1007/s10750-007-9108-z}{Habitat characteristics
of crayfish ({\emph{Paranephrops planifrons}}) in {New Zealand} streams
using generalised additive models ({GAMs})}. Hydrobiologia 596:353--365.

\bibitem[\citeproctext]{ref-Jowett1990}
Jowett IG, Richardson J (1990)
\href{https://doi.org/10.1080/00288330.1990.9516399}{Microhabitat
preferences of benthic invertebrates in a {New Zealand} river and the
development of in-stream flow-habitat models for {\emph{Deleatidium}}
spp.} New Zealand Journal of Marine and Freshwater Research 24:19--30.

\bibitem[\citeproctext]{ref-Joy2013}
Joy M, David B, Lake M (2013) New {Zealand} freshwater fish sampling
protocols: Part 1 --- wadeable rivers and streams. Massey University,
Palmerston North

\bibitem[\citeproctext]{ref-Kusabs2023}
Kusabs IAK (2023) {Lake Rotorua} alum dosing monitoring programme 2023
koura and kakahi. Kusabs \& Associates Ltd, Rotorua, New Zealand

\bibitem[\citeproctext]{ref-Kusabs2026}
Kusabs IAK, Duggan IC, Bruere AC, Scholes P, Te Kurapa TW, Grant TM,
Burdon FJ (2026)
\href{https://doi.org/10.1007/s10530-026-03852-0}{Invasive catfish
linked to the decline of freshwater crayfish in two {Aotearoa New
Zealand} lakes}. Biological Invasions 28.

\bibitem[\citeproctext]{ref-Kusabs2015a}
Kusabs IAK, Hicks BJ, Quinn JM, Hamilton DP (2015)
\href{https://doi.org/10.1016/j.fishres.2015.03.015}{Sustainable
management of freshwater crayfish (kōura, {\emph{Paranephrops
planifrons}}) in {Te Arawa} ({Rotorua}) lakes, {North Island}, {New
Zealand}}. Fisheries Research 168:35--46.

\bibitem[\citeproctext]{ref-Kusabs2009}
Kusabs IAK, Quinn JM (2009)
\href{https://doi.org/10.1080/00288330909510036}{Use of a traditional
{Māori} harvesting method, the tau k{ō}ura, for monitoring k{ō}ura
(freshwater crayfish, {\emph{Paranephrops planifrons}}) in {Lake
Rotoiti}, {North Island, New Zealand}}. New Zealand Journal of Marine
and Freshwater Research 43:713--722.

\bibitem[\citeproctext]{ref-Kusabs2015b}
Kusabs IAK, Quinn JM, Hamilton DP (2015)
\href{https://doi.org/10.1071/mf14148}{Effects of benthic substrate,
nutrient enrichment and predatory fish on freshwater crayfish (koura,
{\emph{Paranephrops planifrons}}) population characteristics in seven
{Te Arawa} ({Rotorua}) lakes, {North Island}, {New Zealand}}. Marine and
Freshwater Research 66:631--643.

\bibitem[\citeproctext]{ref-Kusabs2015c}
Kusabs IAK, Taiaroa R (2015)
\href{https://doi.org/10.13140/RG.2.2.12471.89761}{K{ō}ura populations
in {Lake Taup{ō}} and comments on the effects of catfish}. Ngāti
Tūwharetoa Māori Trust Board

\bibitem[\citeproctext]{ref-Lee2025}
Lee F, Kusabs IAK, Perry GLW, MacNeil C (2025)
\href{https://doi.org/10.5194/we-25-221-2025}{Identifying refugia from
the synergistic threats of climate change and invasive species}. Web
Ecology 25:221--239.

\bibitem[\citeproctext]{ref-Lennox2025}
Lennox RJ, Kambestad M, Berhe S, Birnie‐Gauvin K, Cooke SJ, Dahlmo LS,
Eldøy SH, Davidsen JG, Hanssen EM, Sortland LK, et al. (2025)
\href{https://doi.org/10.1111/rec.14354}{The role of habitat in
predator--prey dynamics with applications to restoration}. Restoration
Ecology 33.

\bibitem[\citeproctext]{ref-Lenth2017}
Lenth RV, Piaskowski J (2017)
\href{https://doi.org/10.32614/cran.package.emmeans}{Emmeans: Estimated
marginal means, aka least-squares means}. CRAN: Contributed Packages.

\bibitem[\citeproctext]{ref-MacNeil2025}
MacNeil C, Lee F, Kusabs IAK, Holmes R (2025)
\href{https://doi.org/10.1051/kmae/2025022}{The biter bit; scavenging
behaviour of native freshwater crayfish on the carrion of native and
introduced fish predators in {Aotearoa-New Zealand}}. Knowledge \&
Management of Aquatic Ecosystems 29.

\bibitem[\citeproctext]{ref-Minns1996}
Minns CK, Kelso JR, Randall RG (1996)
\href{https://doi.org/10.1139/f95-262}{Detecting the response of fish to
habitat alterations in freshwater ecosystems}. Canadian Journal of
Fisheries and Aquatic Sciences 53:403--414.

\bibitem[\citeproctext]{ref-Morelli2016}
Morelli TL, Daly C, Dobrowski SZ, Dulen DM, Ebersole JL, Jackson ST,
Lundquist JD, Millar CI, Maher SP, Monahan WB, et al. (2016)
\href{https://doi.org/10.1371/journal.pone.0159909}{Managing climate
change refugia for climate adaptation}. Rebelo, H (ed.) PLOS ONE
11:e0159909.

\bibitem[\citeproctext]{ref-Ngwenya2025}
Ngwenya DK, Geerts S, Holmes PM, Esler KJ (2025)
\href{https://doi.org/10.1111/rec.70266}{Conceptualizing and
contextualizing {`large‐scale'} and {`scaling‐up'} ecological
restoration}. Restoration Ecology 34.

\bibitem[\citeproctext]{ref-Nystrom2006}
Nyström P, Stenroth P, Holmqvist N, Berglund O, Larsson P, Granéli W
(2006) \href{https://doi.org/10.1111/j.1365-2427.2006.01641.x}{Crayfish
in lakes and streams: Individual and population responses to predation,
productivity and substratum availability}. Freshwater Biology
51:2096--2113.

\bibitem[\citeproctext]{ref-Parkyn2004}
Parkyn SM, Collier KJ (2004)
\href{https://doi.org/10.1023/b:hydr.0000043189.91134.94}{Interaction of
{Press} and {Pulse} disturbance on crayfish populations: Flood impacts
in pasture and forest streams}. Hydrobiologia 527:113--124.

\bibitem[\citeproctext]{ref-Parkyn2001}
Parkyn SM, Collier KJ, Hicks BJ (2001)
\href{https://doi.org/10.1046/j.1365-2427.2001.00702.x}{{New Zealand}
stream crayfish: Functional omnivores but trophic predators?} Freshwater
Biology 46:641--652.

\bibitem[\citeproctext]{ref-Parkyn2009}
Parkyn SM, Meleason MA, Davies-Colley RJ (2009)
\href{https://doi.org/10.1080/00288330909510034}{Wood enhances crayfish
({\emph{Paranephrops planifrons}}) habitat in a forested stream}. New
Zealand Journal of Marine and Freshwater Research 43:689--700.

\bibitem[\citeproctext]{ref-Pilotto2018}
Pilotto F, Nilsson C, Polvi LE, McKie BG (2018)
\href{https://doi.org/10.1002/eap.1672}{First signs of macroinvertebrate
recovery following enhanced restoration of boreal streams used for
timber floating}. Ecological Applications 28:587--597.

\bibitem[\citeproctext]{ref-Prentice2026}
Prentice MJ, Allan MG, Bruere A, Özkundakci D (2026)
\href{https://doi.org/10.1016/j.ecoleng.2025.107862}{Efficacy of an
inflow diversion wall for managing eutrophication in two hydrologically
connected lakes: A modelling study}. Ecological Engineering 223:107862.

\bibitem[\citeproctext]{ref-Quinn2002}
Quinn GP, Keough MJ (2002) Experimental design and data analysis for
biologists. Cambridge University Press

\bibitem[\citeproctext]{ref-RCoreTeam2026}
R Core Team (2026) \href{https://www.R-project.org/}{{R}: A language and
environment for statistical computing}. R Foundation for Statistical
Computing, Vienna, Austria

\bibitem[\citeproctext]{ref-Raven2026b}
Raven OV, Burdon FJ, Kusabs IAK, Holmes R, Özkundakci D (2026)
\href{https://doi.org/10.1007/s10750-026-06302-z}{Littoral habitat
structure influences freshwater crayfish populations in five volcanic
lakes of {Aotearoa New Zealand}}. Hydrobiologia.

\bibitem[\citeproctext]{ref-Raven2026a}
Raven OV, MacNeil C, Lee F, Holmes R, Kusabs IAK, Burdon FJ, Özkundakci
D (2026) \href{https://doi.org/10.21203/rs.3.rs-9891519/v1}{Fear of the
new: A predatory invader provokes stronger non-consumptive effects on a
keystone native prey than a native predator}.

\bibitem[\citeproctext]{ref-Richman2015}
Richman NI, Böhm M, Adams SB, Alvarez F, Bergey EA, Bunn JJS, Burnham Q,
Cordeiro J, Coughran J, Crandall KA, et al. (2015)
\href{https://doi.org/10.1098/rstb.2014.0060}{Multiple drivers of
decline in the global status of freshwater crayfish ({Decapoda}:
{Astacidea})}. Philosophical Transactions of the Royal Society B:
Biological Sciences 370:20140060.

\bibitem[\citeproctext]{ref-Ruh2012}
Ruhí A, Herrmann J, Gascón S, Sala J, Geijer J, Boix D (2012)
\href{https://doi.org/10.1899/11-018.1}{Change in biological traits and
community structure of macroinvertebrates through primary succession in
a man-made {Swedish} wetland}. Freshwater Science 31:22--37.

\bibitem[\citeproctext]{ref-Sayer2025}
Sayer CA, Fernando E, Jimenez RR, Macfarlane NBW, Rapacciuolo G, Böhm M,
Brooks TM, Contreras-MacBeath T, Cox NA, Harrison I, et al. (2025)
\href{https://doi.org/10.1038/s41586-024-08375-z}{One-quarter of
freshwater fauna threatened with extinction}. Nature 638:138--145.

\bibitem[\citeproctext]{ref-Selwood2020}
Selwood KE, Zimmer HC (2020)
\href{https://doi.org/10.1016/j.biocon.2020.108502}{Refuges for
biodiversity conservation: A review of the evidence}. Biological
Conservation 245:108502.

\bibitem[\citeproctext]{ref-Streissl2002}
Streissl F, Hödl W (2002)
\href{https://doi.org/10.1023/A:1021094309738}{Habitat and shelter
requirements of the stone crayfish, {\emph{Austropotamobius torrentium}}
schrank}. Hydrobiologia 477:195--199.

\bibitem[\citeproctext]{ref-Testa2011}
Testa S, Collier KJ, Death RG (2011)
\href{https://doi.org/10.1002/eco.146}{Macroinvertebrate response to
stream restoration by large wood addition}. Ecohydrology 4:800--810.

\bibitem[\citeproctext]{ref-Underwood1994}
Underwood AJ (1994) \href{https://doi.org/10.2307/1942110}{On beyond
{BACI}: Sampling designs that might reliably detect environmental
disturbances}. Ecological Applications 4:3--15.

\bibitem[\citeproctext]{ref-Usio2000}
Usio N, Townsend CR (2000)
\href{https://doi.org/10.1080/00288330.2000.9516957}{Distribution of the
{New Zealand} crayfish ({\emph{Paranephrops zealandicus}}) in relation
to stream physico-chemistry, predatory fish, and invertebrate prey}. New
Zealand Journal of Marine and Freshwater Research 34:557--567.

\bibitem[\citeproctext]{ref-Vehtari2021}
Vehtari A, Gelman A, Simpson D, Carpenter B, Bürkner P-C (2021)
\href{https://doi.org/10.1214/20-ba1221}{Rank-normalization, folding,
and localization: An improved Rˆ for assessing convergence of MCMC (with
discussion)}. Bayesian Analysis 16.

\bibitem[\citeproctext]{ref-Verburg2019}
Verburg P, Albert A (2019) {Lake Taupo} long-term monitoring programme
2017--2018. National Institute of Water \& Atmospheric Research Ltd,
Hamilton, New Zealand

\bibitem[\citeproctext]{ref-Vincent1984}
Vincent WF, Gibbs MM, Dryden SJ (1984)
\href{https://doi.org/10.1080/00288330.1984.9516064}{Accelerated
eutrophication in a {New Zealand} lake: {Lake Rotoiti}, central {North
Island}}. New Zealand Journal of Marine and Freshwater Research
18:431--440.

\bibitem[\citeproctext]{ref-Wentworth1922}
Wentworth CK (1922) \href{https://doi.org/10.1086/622910}{A scale of
grade and class terms for clastic sediments}. The Journal of Geology
30:377--392.

\bibitem[\citeproctext]{ref-Wetzel2001}
Wetzel RG (2001) Limnology: Lake and river ecosystems. 3rd edn. Academic
Press, San Diego, CA

\bibitem[\citeproctext]{ref-Woolway2021}
Woolway RI, Sharma S, Weyhenmeyer GA, Debolskiy A, Golub M,
Mercado-Bettín D, Perroud M, Stepanenko V, Tan Z, Grant L, et al. (2021)
\href{https://doi.org/10.1038/s41467-021-22657-4}{Phenological shifts in
lake stratification under climate change}. Nature Communications 12.

\bibitem[\citeproctext]{ref-Zhao2024}
Zhao M, Feng G, Wang H, Shen C, Fu Y, Zhang Y, Zhang H, Yao Y, Chen J,
Xu W (2024) \href{https://doi.org/10.3390/ani14081147}{The influence of
shelter type and coverage on crayfish ({\emph{Procambarus clarkii}})
predation by catfish ({\emph{Silurus asotus}}): A controlled environment
study}. Animals 14:1147.

\end{CSLReferences}




\end{document}
